\chapter{Human Papillomavirus (HPV)}

\section*{What Is This Test?}
Human papillomavirus (HPV) testing detects high-risk (oncogenic) HPV genotypes in cervical, vaginal, anal, or oropharyngeal specimens. HPV is a non-enveloped, double-stranded circular DNA virus of the \textit{Papillomaviridae} family with >200 identified genotypes, of which approximately 40 infect the anogenital tract. HPV genotypes are classified into high-risk (HR-HPV: 16, 18, 31, 33, 35, 39, 45, 51, 52, 56, 58, 59, 66, 68) and low-risk (LR-HPV: 6, 11 cause >90\% of anogenital warts/condyloma acuminata; others cause cutaneous warts). Persistent infection with HR-HPV is the necessary cause of cervical cancer and its precursors (cervical intraepithelial neoplasia/CIN; squamous intraepithelial lesions/SIL). HPV 16 and 18 together account for 70\% of cervical cancers and the majority of HPV-associated anal, oropharyngeal, vulvar, vaginal, and penile cancers. Most HPV infections are transient; 90\% clear within 2 years. The minority that persist and integrate into the host genome drive progression to high-grade dysplasia and invasive cancer over years to decades. HPV testing is performed either as a standalone test (primary HPV screening) or co-testing with cervical cytology (Pap test) \cite{schiffman2016carcinogenic}.

\section*{Alternative Names}
\begin{itemize}
  \item HPV DNA Test
  \item High-Risk HPV Test
  \item HPV Genotyping
  \item Oncogenic HPV
  \item HPV Co-Testing
  \item Cervical Cancer Screening
  \item HPV Primary Screening
  \item HPV mRNA Test (Aptima)
  \item Cobas HPV Test
\end{itemize}
\section*{Principle}
Multiple FDA-approved assays detect HR-HPV DNA or mRNA: Roche cobas HPV test: Multiplex real-time PCR that detects 14 HR-HPV genotypes with specific genotyping for HPV 16 and HPV 18 (via separate fluorescent channels), plus a pooled result for the other 12 HR-HPV types. The cobas assay is FDA-approved for primary HPV screening (without concomitant cytology) in women 25 and older. Hologic Aptima HPV assay: Transcription-mediated amplification (TMA) targeting HPV E6/E7 mRNA (mRNA markers of oncogene activity rather than viral DNA, theoretically more specific for transforming infections). Detects 14 HR-HPV types but with specific genotyping for HPV 16 and 18/45. BD Onclarity HPV assay: Real-time PCR with extended genotyping (HPV 16, 18, 31, 45, 51, 52, and pooled other HR types). Qiagen Hybrid Capture 2 (HC2): Signal-amplified hybrid capture DNA-RNA hybridization detecting 13 HR-HPV types (older technology, now increasingly replaced by PCR-based tests). ThinPrep and SurePath are the two FDA-approved liquid-based cytology media from which HPV testing can be performed as a reflex from the same vial. Screening algorithms: (1) Primary HPV screening (cobas) starting at age 25. If HPV 16/18+, proceed directly to colposcopy. If positive for other 12 HR-HPV types, reflex cytology; if cytology $\geq$ ASC-US, colposcopy. (2) Co-testing (HPV + cytology) every 5 years starting at age 30. (3) Cytology alone every 3 years starting at age 21.

\section*{History}
The causal link between HPV and cervical cancer was established by Harald zur Hausen in the 1980s (Nobel Prize 2008). The first HPV DNA test (Hybrid Capture) was FDA-approved in 1999. Co-testing (HPV + cytology) gained FDA approval in 2003 and was incorporated into screening guidelines. The cobas HPV test was FDA-approved for primary screening in 2014, and the ASCCP guidelines endorsed primary HPV screening in 2015 \cite{huh2015use}. The HPV vaccines (Gardasil, quadrivalent HPV 6/11/16/18, 2006; Cervarix, bivalent HPV 16/18, 2009; Gardasil 9, 9-valent HPV 6/11/16/18/31/33/45/52/58, 2014) are the primary prevention strategy, with vaccination recommended for all girls and boys at age 11--12 (can start at age 9) and catch-up vaccination through age 26 (shared clinical decision-making for ages 27--45) \cite{markowitz2014human}. The ASCCP Risk-Based Management Consensus Guidelines (2019) introduced risk-based management thresholds for colposcopy and treatment. Widespread HPV vaccination has dramatically reduced HPV prevalence, genital warts, and cervical precancer in vaccinated cohorts.

\section*{Why This Test Is Ordered}
\begin{itemize}
  \item Primary cervical cancer screening: follow the current age-, risk-, and screening-history-specific national guideline; HPV testing intervals and eligible ages vary by guideline and assay authorization \cite{uspstf2018cervical}
  \item Triage of abnormal cervical cytology results: HPV testing is recommended for women with ASC-US (atypical squamous cells of undetermined significance) cytology to determine need for colposcopy. If HR-HPV is positive, proceed to colposcopy; if HR-HPV is negative, return to routine screening
  \item Post-treatment surveillance: After treatment for CIN 2/3 (LEEP, conization), co-testing at 12 and 24 months is recommended to document clearance of HR-HPV and ensure no residual/recurrent disease. If both tests are negative at both time points, the patient may return to extended-interval screening
  \item Evaluation of women with persistent HR-HPV positivity (positive HPV at two consecutive screenings over 1--2 years) with normal or ASC-US cytology; increased risk of underlying CIN 3+ warrants colposcopy
  \item Co-testing for women over 65 to determine eligibility to exit screening: If the patient has had adequate negative screening in the past 10 years (3 consecutive negative cytology tests or 2 consecutive negative co-tests, with the most recent within 5 years) and no history of CIN 2+, screening may be discontinued
  \item Anal cancer screening (high-risk populations): Anal HPV with or without anal cytology in HIV-infected MSM, women with history of cervical/vulvar dysplasia, and solid organ transplant recipients. However, standardized screening guidelines are less established than for cervical cancer
\end{itemize}

\section*{Primary vs Secondary Test}
Primary HPV testing (without cytology) is FDA-approved as a primary cervical cancer screening test for women 25 and older. HPV testing is also a secondary triage test for ASC-US cytology and for post-colposcopy surveillance to risk-stratify management. Co-testing (HPV + cytology) places both tests as primary. HPV testing is NOT used as a primary screening test for women under 25 (due to high prevalence of transient HPV infection and low risk of cancer).

\section*{How This Test Is Performed}
Cervical specimen collection: A cervical sample is collected using a broom-style device or endocervical brush plus spatula (similar to a Pap smear). The collection device is rotated in the endocervical canal (10--15 rotations) and then rinsed vigorously into the liquid-based cytology vial (ThinPrep or SurePath preservative medium). The same vial serves for both cytology and HPV testing (reflex HPV testing from residual material). Vaginal self-collection for HPV testing is currently under investigation and used in some countries but is not yet FDA-approved in the United States for cervical cancer screening (except for use in clinical trials and some limited clinical settings). Adequate cellularity of the specimen is required for valid results.

\section*{What Preparation Is Needed}
The patient should avoid douching, using vaginal creams, suppositories, or spermicides for 48 hours before specimen collection. Sexual intercourse and use of tampons should be avoided for 24--48 hours before specimen collection. The specimen should ideally not be collected during heavy menstruation, as blood may obscure cytology interpretation and potentially interfere with HPV testing. However, HPV testing alone is more tolerant of menstrual blood than cytology.

\section*{Contraindications and Precautions}
\begin{itemize}
  \item HPV testing is NOT recommended in women under 21 years (no screening at all) or in women aged 21--24 as a primary screening method because the prevalence of transient HR-HPV infection is extremely high (up to 40--60\%), and the risk of cervical cancer in this age group is near zero. HPV testing in this age group leads to excessive colposcopy referrals without reducing cancer incidence. HPV testing is NOT recommended for screening in women who have undergone total hysterectomy (removal of the uterus and cervix) for benign disease, because there is no cervix to screen. Women with a history of CIN 2+ treated by hysterectomy should continue vaginal cuff cytology and HPV co-testing per guidelines. A positive HR-HPV test does NOT mean the patient has cancer, nor does it necessarily mean she has CIN 2/3. It indicates the presence of a high-risk HPV infection, which in the vast majority of cases will clear without intervention. Colposcopy should be performed based on the ASCCP risk-based management algorithm (immediate risk of CIN 3+ >4\% warrants colposcopy). A negative cytology in the setting of a positive HPV test carries a much lower risk of CIN 3+ than a positive cytology.
\end{itemize}
\section*{Risks}
Cervical specimen collection is minimally invasive. The primary risks of HPV testing pertain to psychological and behavioral consequences: anxiety associated with a positive HPV result (due to the association with "cancer" and "STI"), stigma, and concerns about disclosure. Patient education is essential: a positive HPV test is common, most infections clear, and the purpose of screening is to identify persistent infections to prevent cervical cancer, not to discover disease. Over-screening and over-treatment of transient HPV infections are the primary system-level risks.

\section*{What Happens After the Test}
HPV results available in 1--7 days. Management is determined by the ASCCP Risk-Based Management Consensus Guidelines (2019), based on the immediate risk of CIN 3+ \cite{perkins2020asccp}:
  - HPV 16 or 18 positive (primary HPV screening): Proceed directly to colposcopy regardless of cytology result (immediate CIN 3+ risk >4\%)
  - HPV positive (other 12 HR-HPV types) with reflex cytology negative: Repeat co-testing in 1 year
  - HPV positive with cytology ASC-US or greater: Colposcopy
  - HPV negative with cytology negative: Return at 5-year interval
  - Co-testing both negative: 5-year interval
  - HPV negative, cytology ASC-US: Repeat co-testing in 3 years (low risk)

\section*{Understanding Results}

\subsection*{Below Normal}
\begin{itemize}
  \item \textbf{HR-HPV not detected (negative):} No high-risk HPV DNA (or mRNA) detected. In a woman with negative cytology, this result provides the highest level of reassurance (5-year CIN 3+ risk <0.15\%). Routine screening interval (5 years for primary HPV or co-testing; 3 years for cytology alone) applies. In a woman with ASC-US cytology, a negative HPV test indicates that the ASC-US is benign and she should return to routine screening in 3 years (no colposcopy)
  \item \textbf{Cytology negative for intraepithelial lesion or malignancy (NILM) with negative HPV:} Return to screening in 5 years (co-testing) or at the interval indicated by the screening strategy used
\end{itemize}

\subsection*{Above Normal}
\begin{itemize}
  \item \textbf{HR-HPV positive (detected):} High-risk HPV nucleic acid detected. Does NOT indicate cancer; does NOT necessarily indicate precancer. Indicates the presence of a high-risk HPV infection. A single positive test should be managed according to the ASCCP algorithm based on the specific HPV result and cytology. Persistent detection (positive at two consecutive screens 12 months apart) is the biomarker for increased cancer risk and warrants colposcopy
  \item \textbf{HPV 16 and/or 18 positive:} Highest-risk HPV genotypes. Colposcopy is recommended regardless of cytology result, even if cytology is NILM. HPV 16 has the strongest association with invasive cervical squamous cell carcinoma; HPV 18 is proportionally overrepresented in adenocarcinoma
  \item \textbf{HPV positive, other 12 HR-HPV types (HPV 16/18 negative):} Risk is intermediate. Next step depends on cytology: if cytology is NILM, repeat co-testing in 1 year; if cytology $\geq$ ASC-US, colposcopy
  \item \textbf{Co-testing: HPV positive, cytology abnormal (ASC-US or greater):} Colposcopy. The combination of positive HPV and abnormal cytology (especially HSIL or ASC-H) confers the highest immediate risk of CIN 3+ (>60\% for HSIL; >25\% for ASC-H)
  \item \textbf{Cytology HSIL (high-grade squamous intraepithelial lesion):} Immediate colposcopy regardless of HPV result; HPV testing is not required for triage because HSIL carries a 60--80\% risk of CIN 2+ and a 2\% risk of invasive cancer
\end{itemize}

\section*{Normal Results}
\textbf{High-risk HPV not detected (negative)}. For women 30--65, co-testing: \textbf{HPV not detected and cytology NILM (negative for intraepithelial lesion or malignancy)} is the optimal screening result.

\section*{Follow-up Tests}
\begin{itemize}
  \item \textbf{Colposcopy:} Indicated based on ASCCP risk thresholds (immediate CIN 3+ risk $\geq$4\%). Colposcopy with application of 3--5\% acetic acid identifies acetowhite lesions; biopsies are taken of any visible lesions. Endocervical curettage (ECC) may be performed
  \item \textbf{Repeat co-testing (HPV + cytology) in 1 year:} For patients with HPV positive (non-16/18) and NILM cytology, this is the recommended surveillance interval. If both tests are negative at 12 months, return to routine screening interval
  \item \textbf{Cervical biopsy (colposcopy-directed and/or random):} Histopathologic diagnosis (CIN 1, 2, 3, or invasive carcinoma). CIN 1 is managed expectantly (likely to regress); CIN 2/3 is treated by excision (LEEP) or ablation
  \item \textbf{HPV genotyping:} For risk stratification when not already performed; identification of HPV 16/18 can alter management
  \item \textbf{p16 immunostaining (on cervical biopsy):} An adjunctive histopathologic stain that helps distinguish CIN 2/3 (which overexpresses p16) from benign mimics (immature squamous metaplasia, atrophy); improves diagnostic reproducibility
  \item \textbf{Endocervical curettage (ECC):} To evaluate the endocervical canal when the squamocolumnar junction is not fully visualized, or when cytology suggests glandular abnormality (AGC---atypical glandular cells)
\end{itemize}

\section*{References}
\bibliographystyle{unsrtnat}
\bibliography{references}

\clearpage
\section*{Regulatory Status and Authorization Date}
This chapter describes a test category, not a specific commercial product. FDA, CDSCO, EU IVDR, and MHRA approval or registration dates therefore cannot be assigned without the assay manufacturer, product name, instrument, and intended use. For laboratory-developed tests, an individual agency approval date may not exist. See the Preface for the interpretation of regulatory status.

\clearpage
